% ============================================================
%	附录：代数基础各定理与主要性质的证明
%	本附录给出第 2 章代数基础中定理与主要性质的证明。
% ============================================================

\section{谱分解定理的证明}
\label{sec:谱分解证明}

\begin{proof}
设 $\boldsymbol{A}$ 为 $n$ 阶实对称矩阵，分三步证明式~\eqref{eq:矩阵:谱分解}。

第一步，证明特征值全为实数。设 $\lambda$ 为 $\boldsymbol{A}$ 的特征值，$\boldsymbol{v}$ 为对应的单位特征向量（$\boldsymbol{v}^{\mathrm{H}}\boldsymbol{v} = 1$），则 $\boldsymbol{A}\boldsymbol{v} = \lambda\boldsymbol{v}$。左乘 $\boldsymbol{v}^{\mathrm{H}}$ 得 $\lambda = \boldsymbol{v}^{\mathrm{H}}\boldsymbol{A}\boldsymbol{v}$，对实对称矩阵 $\boldsymbol{v}^{\mathrm{H}}\boldsymbol{A}\boldsymbol{v} = \boldsymbol{v}^{\mathrm{T}}\boldsymbol{A}\boldsymbol{v}\in\mathbb{R}$；又 $\boldsymbol{A} = \boldsymbol{A}^{\mathrm{T}} = \boldsymbol{A}^{\mathrm{H}}$，故
\[
\lambda = \boldsymbol{v}^{\mathrm{H}}\boldsymbol{A}\boldsymbol{v} = (\boldsymbol{v}^{\mathrm{H}}\boldsymbol{A}\boldsymbol{v})^{\mathrm{H}} = \lambda^{*},
\]
即 $\lambda = \lambda^{*}$，特征值为实数。

第二步，证明不同特征值的特征向量相互正交。设 $\lambda_1\neq\lambda_2$，对应的特征向量为 $\boldsymbol{v}_1$、$\boldsymbol{v}_2$，则
\[
\lambda_1\boldsymbol{v}_2^{\mathrm{T}}\boldsymbol{v}_1 = \boldsymbol{v}_2^{\mathrm{T}}\boldsymbol{A}\boldsymbol{v}_1 = (\boldsymbol{A}\boldsymbol{v}_2)^{\mathrm{T}}\boldsymbol{v}_1 = \lambda_2\boldsymbol{v}_2^{\mathrm{T}}\boldsymbol{v}_1,
\]
故 $(\lambda_1-\lambda_2)\boldsymbol{v}_2^{\mathrm{T}}\boldsymbol{v}_1 = 0$，由 $\lambda_1\neq\lambda_2$ 得 $\boldsymbol{v}_1\perp\boldsymbol{v}_2$。

第三步，构造正交矩阵。对特征值作归纳：$n=1$ 时结论显然。设 $n-1$ 阶结论成立。取单位特征向量 $\boldsymbol{v}_1$，补充与 $\boldsymbol{v}_1$ 正交的单位向量 $\boldsymbol{v}_2,\cdots,\boldsymbol{v}_n$，令 $\boldsymbol{Q}_1 = [\boldsymbol{v}_1\ \boldsymbol{v}_2\ \cdots\ \boldsymbol{v}_n]$ 为正交矩阵，则
\[
\boldsymbol{Q}_1^{\mathrm{T}}\boldsymbol{A}\boldsymbol{Q}_1 = \begin{bmatrix} \lambda_1 & \boldsymbol{0} \\ \boldsymbol{0} & \boldsymbol{A}_1 \end{bmatrix},
\]
式中 $\boldsymbol{A}_1$ 仍为实对称矩阵（$(\boldsymbol{Q}_1^{\mathrm{T}}\boldsymbol{A}\boldsymbol{Q}_1)^{\mathrm{T}} = \boldsymbol{Q}_1^{\mathrm{T}}\boldsymbol{A}\boldsymbol{Q}_1$）。由归纳假设存在正交矩阵 $\boldsymbol{Q}_2$ 使 $\boldsymbol{Q}_2^{\mathrm{T}}\boldsymbol{A}_1\boldsymbol{Q}_2$ 对角，于是 $\boldsymbol{Q} = \boldsymbol{Q}_1\operatorname{diag}(1,\boldsymbol{Q}_2)$ 即满足式~\eqref{eq:矩阵:谱分解}。证毕。
\end{proof}

\section{Jordan 标准形定理的证明}
\label{sec:jordan证明}

\begin{proof}
分四步证明式~\eqref{eq:矩阵:jordan标准形}。

第一步（Schur 三角化）：对特征值作归纳可证，任意复方阵 $\boldsymbol{A}$ 相似于上三角矩阵，且主对角线元素即其特征值（重数计入）。故不妨设 $\boldsymbol{A}$ 已是上三角矩阵，特征值沿主对角线排列。

第二步（按特征值分块）：把 $\boldsymbol{A}$ 的主对角线按相同的特征值分组，同时把相邻行与列同时交换（相似变换），得块上三角矩阵，各对角块 $\boldsymbol{A}_i$ 的对角线元素全为同一 $\lambda_i$。

第三步（对角块化为 Jordan 块）：考察零对角线矩阵 $\boldsymbol{N}_i = \boldsymbol{A}_i-\lambda_i\boldsymbol{I}$（幂零）。对幂零矩阵存在（唯一的）按广义特征空间的循环分解：值域链
\[
\operatorname{Im}\,\boldsymbol{N}_i^{k} \subsetneq \cdots \subsetneq \operatorname{Im}\,\boldsymbol{N}_i \subsetneq \mathbb{C}^{k_i}
\]
给出基，在该基下 $\boldsymbol{N}_i$ 化为由若干形如式~\eqref{eq:矩阵:jordan块} 的 $0$ 对角 Jordan 块（各块主对角线上方为 1）构成的块对角矩阵。相应的 $\boldsymbol{A}_i = \lambda_i\boldsymbol{I}+\boldsymbol{N}_i$ 化为若干个 $\boldsymbol{J}_{k}(\lambda_i)$。

第四步（回代与唯一性）：把各步的相似变换矩阵相乘即得 $\boldsymbol{P}$。特征值 $\lambda_i$ 的代数重数等于相应 Jordan 块阶数之和，几何重数等于 Jordan 块个数，二者均由 $\boldsymbol{A}$ 的秩与特征子空间链唯一决定，故除 Jordan 块排列顺序外，标准形唯一。证毕。
\end{proof}

\section{Woodbury 恒等式的证明}
\label{sec:woodbury证明}

\begin{proof}
只需验证式~\eqref{eq:矩阵:woodbury} 右端为 $(\boldsymbol{A}+\boldsymbol{U}\boldsymbol{B}\boldsymbol{V})$ 的逆。记
\[
\boldsymbol{S} = \boldsymbol{B}^{-1}+\boldsymbol{V}\boldsymbol{A}^{-1}\boldsymbol{U}, \qquad
\boldsymbol{M} = \boldsymbol{A}+\boldsymbol{U}\boldsymbol{B}\boldsymbol{V},
\]
则
\begin{align*}
&\big[\boldsymbol{A}^{-1}-\boldsymbol{A}^{-1}\boldsymbol{U}\boldsymbol{S}^{-1}\boldsymbol{V}\boldsymbol{A}^{-1}\big]\,\boldsymbol{M} \\
&= \boldsymbol{I}+\boldsymbol{A}^{-1}\boldsymbol{U}\boldsymbol{B}\boldsymbol{V}-\boldsymbol{A}^{-1}\boldsymbol{U}\boldsymbol{S}^{-1}\boldsymbol{V}-\boldsymbol{A}^{-1}\boldsymbol{U}\boldsymbol{S}^{-1}\boldsymbol{V}\boldsymbol{A}^{-1}\boldsymbol{U}\boldsymbol{B}\boldsymbol{V} \\
&= \boldsymbol{I}+\boldsymbol{A}^{-1}\boldsymbol{U}\big[\boldsymbol{B}\boldsymbol{V}-\boldsymbol{S}^{-1}(\boldsymbol{S}\boldsymbol{B}\boldsymbol{V}-\boldsymbol{V})\big] \\
&= \boldsymbol{I}+\boldsymbol{A}^{-1}\boldsymbol{U}\big[\boldsymbol{B}\boldsymbol{V}-\boldsymbol{S}^{-1}(\boldsymbol{B}^{-1}\boldsymbol{B}\boldsymbol{V}+\boldsymbol{V}\boldsymbol{A}^{-1}\boldsymbol{U}\boldsymbol{B}\boldsymbol{V}-\boldsymbol{V})\big] \\
&= \boldsymbol{I}+\boldsymbol{A}^{-1}\boldsymbol{U}\big[\boldsymbol{B}\boldsymbol{V}-\boldsymbol{S}^{-1}\boldsymbol{V}\boldsymbol{A}^{-1}\boldsymbol{U}\boldsymbol{B}\boldsymbol{V}\big] \\
&= \boldsymbol{I}+\boldsymbol{A}^{-1}\boldsymbol{U}\boldsymbol{S}^{-1}\big[\boldsymbol{S}\boldsymbol{B}\boldsymbol{V}-\boldsymbol{V}\boldsymbol{A}^{-1}\boldsymbol{U}\boldsymbol{B}\boldsymbol{V}\big] \\
&= \boldsymbol{I}+\boldsymbol{A}^{-1}\boldsymbol{U}\boldsymbol{S}^{-1}\big[(\boldsymbol{B}^{-1}+\boldsymbol{V}\boldsymbol{A}^{-1}\boldsymbol{U})\boldsymbol{B}\boldsymbol{V}-\boldsymbol{V}\boldsymbol{A}^{-1}\boldsymbol{U}\boldsymbol{B}\boldsymbol{V}\big] \\
&= \boldsymbol{I}+\boldsymbol{A}^{-1}\boldsymbol{U}\boldsymbol{S}^{-1}\boldsymbol{V} = \boldsymbol{I}+\boldsymbol{A}^{-1}\boldsymbol{U}\boldsymbol{S}^{-1}\boldsymbol{V}-\boldsymbol{A}^{-1}\boldsymbol{U}\boldsymbol{S}^{-1}\boldsymbol{V} = \boldsymbol{I}.
\end{align*}
故式~\eqref{eq:矩阵:woodbury} 右端为 $\boldsymbol{M}^{-1}$。证毕。
\end{proof}

\section{奇异值分解的证明}
\label{sec:SVD证明}

\begin{proof}
设 $\boldsymbol{A}\in\mathbb{R}^{n\times m}$。矩阵 $\boldsymbol{A}^{\mathrm{T}}\boldsymbol{A}\in\mathbb{R}^{m\times m}$ 实对称半正定，由谱分解（式~\eqref{eq:矩阵:谱分解}）存在正交矩阵 $\boldsymbol{V}\in\mathbb{R}^{m\times m}$ 使
\[
\boldsymbol{V}^{\mathrm{T}}\boldsymbol{A}^{\mathrm{T}}\boldsymbol{A}\boldsymbol{V} = \operatorname{diag}(\sigma_1^{2},\cdots,\sigma_m^{2}),
\]
式中 $\sigma_i^{2}\geq 0$ 为特征值。设 $\sigma_1\geq\cdots\geq\sigma_r>0$、$\sigma_{r+1}=\cdots=\sigma_m=0$，$r = \operatorname{rank}(\boldsymbol{A})$。记 $\boldsymbol{V} = [\boldsymbol{v}_1\ \cdots\ \boldsymbol{v}_m]$，则
\[
\boldsymbol{A}^{\mathrm{T}}\boldsymbol{A}\boldsymbol{v}_i = \sigma_i^{2}\boldsymbol{v}_i,\qquad \|\boldsymbol{A}\boldsymbol{v}_i\|^{2} = \boldsymbol{v}_i^{\mathrm{T}}\boldsymbol{A}^{\mathrm{T}}\boldsymbol{A}\boldsymbol{v}_i = \sigma_i^{2}.
\]
令 $\boldsymbol{u}_i = \boldsymbol{A}\boldsymbol{v}_i/\sigma_i$（$i=1,\cdots,r$），则 $\{\boldsymbol{u}_1,\cdots,\boldsymbol{u}_r\}$ 为正交单位组；补充 $\boldsymbol{u}_{r+1},\cdots,\boldsymbol{u}_n$ 使 $\boldsymbol{U} = [\boldsymbol{u}_1\ \cdots\ \boldsymbol{u}_n]$ 为正交矩阵。于是对 $i=1,\cdots,r$ 有 $\boldsymbol{A}\boldsymbol{v}_i = \sigma_i\boldsymbol{u}_i$，对 $i>r$ 有 $\boldsymbol{A}\boldsymbol{v}_i = \boldsymbol{0}$，即
\[
\boldsymbol{A}\boldsymbol{V} = \boldsymbol{U}\boldsymbol{D},
\]
式中 $\boldsymbol{D}\in\mathbb{R}^{n\times m}$ 的对角元素为 $\sigma_1,\cdots,\sigma_r,0,\cdots$。右乘 $\boldsymbol{V}^{\mathrm{T}}$ 得式~\eqref{eq:矩阵:SVD}。$\boldsymbol{u}_i = \boldsymbol{A}\boldsymbol{v}_i/\sigma_i$ 满足 $\boldsymbol{A}\boldsymbol{A}^{\mathrm{T}}\boldsymbol{u}_i = \sigma_i^{2}\boldsymbol{u}_i$，即 $\boldsymbol{U}$ 的列向量为 $\boldsymbol{A}\boldsymbol{A}^{\mathrm{T}}$ 的单位特征向量。证毕。
\end{proof}

\section{LU 分解与 Cholesky 分解的证明}
\label{sec:LU证明}

\begin{proof}
先证 LU 分解（式~\eqref{eq:矩阵:LU}）。设 $\boldsymbol{A}$ 的各阶顺序主子式非零，对阶数 $n$ 作归纳。$n=1$ 时 $\boldsymbol{A} = [a_{11}] = [1][a_{11}]$。设 $n-1$ 阶成立。把 $\boldsymbol{A}$ 分块为
\[
\boldsymbol{A} = \begin{bmatrix} \boldsymbol{A}_{n-1} & \boldsymbol{a} \\ \boldsymbol{b}^{\mathrm{T}} & a_{nn} \end{bmatrix},
\]
由归纳假设 $\boldsymbol{A}_{n-1} = \boldsymbol{L}_{n-1}\boldsymbol{U}_{n-1}$（可逆）。令
\[
\boldsymbol{L} = \begin{bmatrix} \boldsymbol{L}_{n-1} & \boldsymbol{0} \\ \boldsymbol{b}^{\mathrm{T}}\boldsymbol{U}_{n-1}^{-1} & 1 \end{bmatrix}, \qquad
\boldsymbol{U} = \begin{bmatrix} \boldsymbol{U}_{n-1} & \boldsymbol{L}_{n-1}^{-1}\boldsymbol{a} \\ \boldsymbol{0} & a_{nn}-\boldsymbol{b}^{\mathrm{T}}\boldsymbol{A}_{n-1}^{-1}\boldsymbol{a} \end{bmatrix},
\]
直接相乘即得 $\boldsymbol{A} = \boldsymbol{L}\boldsymbol{U}$，且 $\boldsymbol{L}$ 单位下三角、$\boldsymbol{U}$ 上三角。唯一性由分解两端的下三角/上三角子矩阵逐步比较得到。

再证 Cholesky 分解（式~\eqref{eq:矩阵:cholesky}）。$\boldsymbol{A}$ 对称正定时由 LU 分解 $\boldsymbol{A} = \boldsymbol{L}\boldsymbol{U}$，$\boldsymbol{U} = \boldsymbol{L}^{-1}\boldsymbol{A}$ 亦对称正定，故 $\boldsymbol{U} = \boldsymbol{D}\boldsymbol{L}^{\mathrm{T}}$（$\boldsymbol{D}$ 对角、$\boldsymbol{L}$ 单位下三角），即 $\boldsymbol{A} = \boldsymbol{L}\boldsymbol{D}\boldsymbol{L}^{\mathrm{T}}$（LDL 分解）。由 $\boldsymbol{A}$ 正定，$\boldsymbol{D}$ 对角元为正，记 $\boldsymbol{L}\sqrt{\boldsymbol{D}} = \hat{\boldsymbol{L}}$，则 $\boldsymbol{A} = \hat{\boldsymbol{L}}\hat{\boldsymbol{L}}^{\mathrm{T}}$。证毕。
\end{proof}

\section{特征值性质的证明}
\label{sec:特征值证明}

\subsection*{特征值与迹、行列式的关系}

\begin{proof}
式~\eqref{eq:矩阵:特征值迹行列式} 由特征多项式的因式分解与 Vieta 定理得到。特征多项式（式~\eqref{eq:矩阵:特征多项式}）可写为
\[
\det(\lambda\boldsymbol{I}-\boldsymbol{A}) = \prod_{i=1}^{n}(\lambda-\lambda_i) = \lambda^{n}-(g_1)\lambda^{n-1}+\cdots+(-1)^{n}g_n,
\]
式中 $g_1 = \sum_{i}\lambda_i$、$g_n = \prod_{i}\lambda_i$。另一方面，把 $\det(\lambda\boldsymbol{I}-\boldsymbol{A})$ 按行列式定义展开，含 $\lambda^{n-1}$ 的项恰来自主对角线乘积 $(\lambda-a_{11})\cdots(\lambda-a_{nn})$，故
\[
g_1 = \operatorname{tr}(\boldsymbol{A}),\qquad g_n = \det\boldsymbol{A}.
\]
于是 $\operatorname{tr}(\boldsymbol{A}) = \sum_{i}\lambda_i$、$\det\boldsymbol{A} = \prod_{i}\lambda_i$。证毕。
\end{proof}

\subsection*{AB 与 BA 的非零特征值相同}

\begin{proof}
设 $\lambda\neq 0$ 为 $\boldsymbol{A}\boldsymbol{B}$ 的特征值，$\boldsymbol{v}\neq\boldsymbol{0}$ 满足 $\boldsymbol{A}\boldsymbol{B}\boldsymbol{v} = \lambda\boldsymbol{v}$。则 $\boldsymbol{B}\boldsymbol{v}\neq\boldsymbol{0}$（否则 $\lambda\boldsymbol{v} = \boldsymbol{0}$），且
\[
\boldsymbol{B}\boldsymbol{A}(\boldsymbol{B}\boldsymbol{v}) = \boldsymbol{B}(\boldsymbol{A}\boldsymbol{B}\boldsymbol{v}) = \lambda(\boldsymbol{B}\boldsymbol{v}),
\]
故 $\lambda$ 也是 $\boldsymbol{B}\boldsymbol{A}$ 的特征值。由对称性两者非零特征值完全相同。当 $\lambda = 0$ 时结论需另行处理（代数重数可能不同），但非零特征值集合一致。证毕。
\end{proof}

\section{正定矩阵性质的证明}
\label{sec:正定证明}

\subsection*{特征值判别法}

\begin{proof}
$\boldsymbol{A}$ 实对称时由谱分解（式~\eqref{eq:矩阵:谱分解}）$\boldsymbol{A} = \boldsymbol{Q}\operatorname{diag}(\lambda_i)\boldsymbol{Q}^{\mathrm{T}}$，故对任意 $\boldsymbol{x}\neq\boldsymbol{0}$，记 $\boldsymbol{y} = \boldsymbol{Q}^{\mathrm{T}}\boldsymbol{x}\neq\boldsymbol{0}$，
\[
\boldsymbol{x}^{\mathrm{T}}\boldsymbol{A}\boldsymbol{x} = \sum_{i=1}^{n}\lambda_i y_i^{2}.
\]
由于 $y_i$ 可取任意不全为零的实数组，上式对所有 $\boldsymbol{x}\neq\boldsymbol{0}$ 恒正当且仅当全部 $\lambda_i > 0$。证毕。
\end{proof}

\subsection*{顺序主子式判别法（Sylvester 准则）}

\begin{proof}
必要性：若 $\boldsymbol{A}$ 正定，对前 $k$ 个分量的子向量 $\boldsymbol{x}\neq\boldsymbol{0}$ 有 $\boldsymbol{x}^{\mathrm{T}}\boldsymbol{A}_k\boldsymbol{x} = \boldsymbol{x}^{\mathrm{T}}\boldsymbol{A}\boldsymbol{x} > 0$，故各顺序主子矩阵 $\boldsymbol{A}_k$ 正定，$\det\boldsymbol{A}_k > 0$。

充分性：设全部顺序主子式 $\det\boldsymbol{A}_k > 0$。由 LU 分解（式~\eqref{eq:矩阵:LU}）$\boldsymbol{A} = \boldsymbol{L}\boldsymbol{U}$，且 $\det\boldsymbol{A}_k = \prod_{i=1}^{k}U_{ii}$，故 $U_{ii}>0$。于是 $\boldsymbol{A} = \boldsymbol{L}\boldsymbol{D}\boldsymbol{L}^{\mathrm{T}}$，$\boldsymbol{D}$ 对角元为正；令 $\boldsymbol{B} = \boldsymbol{L}\sqrt{\boldsymbol{D}}$（下三角），则 $\boldsymbol{A} = \boldsymbol{B}\boldsymbol{B}^{\mathrm{T}}$，对任意 $\boldsymbol{x}\neq\boldsymbol{0}$ 有 $\boldsymbol{x}^{\mathrm{T}}\boldsymbol{A}\boldsymbol{x} = \|\boldsymbol{B}^{\mathrm{T}}\boldsymbol{x}\|^{2} > 0$（$\boldsymbol{B}$ 可逆）。故 $\boldsymbol{A}$ 正定。证毕。
\end{proof}

\subsection*{秩不等式 $\operatorname{rank}(\boldsymbol{B}\boldsymbol{A}\boldsymbol{B}^{\mathrm{T}}) = \operatorname{rank}(\boldsymbol{B})$}

\begin{proof}
$\boldsymbol{A}$ 正定故可分解 $\boldsymbol{A} = \boldsymbol{C}\boldsymbol{C}^{\mathrm{T}}$（$\boldsymbol{C}$ 可逆），于是 $\boldsymbol{B}\boldsymbol{A}\boldsymbol{B}^{\mathrm{T}} = (\boldsymbol{B}\boldsymbol{C})(\boldsymbol{B}\boldsymbol{C})^{\mathrm{T}}$。由秩恒等式 $\operatorname{rank}(\boldsymbol{X}\boldsymbol{X}^{\mathrm{T}}) = \operatorname{rank}(\boldsymbol{X})$（见式~\eqref{eq:矩阵:rank恒等式}），
\[
\operatorname{rank}(\boldsymbol{B}\boldsymbol{A}\boldsymbol{B}^{\mathrm{T}}) = \operatorname{rank}(\boldsymbol{B}\boldsymbol{C}) = \operatorname{rank}(\boldsymbol{B}),
\]
末等号因 $\boldsymbol{C}$ 可逆。证毕。
\end{proof}

\subsection*{半正定的 $X^{\mathrm{T}}AX = 0$ 性质}

\begin{proof}
设 $\boldsymbol{A}$ 半正定且 $\boldsymbol{X}^{\mathrm{T}}\boldsymbol{A}\boldsymbol{X} = \boldsymbol{0}$。对任意 $\boldsymbol{y}$，$\boldsymbol{z} = \boldsymbol{X}\boldsymbol{y}$，则 $\boldsymbol{z}^{\mathrm{T}}\boldsymbol{A}\boldsymbol{z} = \boldsymbol{y}^{\mathrm{T}}\boldsymbol{X}^{\mathrm{T}}\boldsymbol{A}\boldsymbol{X}\boldsymbol{y} = 0$。由 $\boldsymbol{A}$ 半正定，$\boldsymbol{z}^{\mathrm{T}}\boldsymbol{A}\boldsymbol{z} = 0$ 蕴含 $\boldsymbol{A}\boldsymbol{z} = \boldsymbol{0}$（记 $\boldsymbol{A} = \boldsymbol{C}\boldsymbol{C}^{\mathrm{T}}$，则 $\boldsymbol{z}^{\mathrm{T}}\boldsymbol{A}\boldsymbol{z} = \|\boldsymbol{C}^{\mathrm{T}}\boldsymbol{z}\|^{2} = 0\Rightarrow \boldsymbol{C}^{\mathrm{T}}\boldsymbol{z} = \boldsymbol{0}\Rightarrow \boldsymbol{A}\boldsymbol{z} = \boldsymbol{0}$）。对一切 $\boldsymbol{y}$ 成立，故 $\boldsymbol{A}\boldsymbol{X} = \boldsymbol{0}$。证毕。
\end{proof}

\section{范数不等式}
\label{sec:范数证明}

\subsection*{Cauchy-Schwarz 不等式}

\begin{proof}
对任意 $\boldsymbol{x},\boldsymbol{y}\in\mathbb{R}^n$ 与实数 $t$，
\[
0 \leq \|\boldsymbol{x}+t\boldsymbol{y}\|_2^{2} = \|\boldsymbol{x}\|_2^{2}+2t\boldsymbol{x}^{\mathrm{T}}\boldsymbol{y}+t^{2}\|\boldsymbol{y}\|_2^{2}.
\]
右端是关于 $t$ 的二次式恒非负，判别式非正：
\[
(2\boldsymbol{x}^{\mathrm{T}}\boldsymbol{y})^{2}-4\|\boldsymbol{x}\|_2^{2}\|\boldsymbol{y}\|_2^{2}\leq 0,
\]
即 $|\boldsymbol{x}^{\mathrm{T}}\boldsymbol{y}|\leq\|\boldsymbol{x}\|_2\|\boldsymbol{y}\|_2$。等号当且仅当判别式为零，即 $\boldsymbol{x}+t\boldsymbol{y} = \boldsymbol{0}$ 有实根，两向量线性相关。证毕。
\end{proof}

\section{伪逆性质与导数的证明}
\label{sec:伪逆证明}

\subsection*{伪逆的性质}

\begin{proof}
由定义（式~\eqref{eq:矩阵:伪逆}）直接验证四个条件可证各性质。
$(\boldsymbol{A}^{+})^{+} = \boldsymbol{A}$：$\boldsymbol{A}^{+}$ 满足条件等价于 $\boldsymbol{A}$ 对 $\boldsymbol{A}^{+}$ 亦满足（四条件在 $\boldsymbol{A}\leftrightarrow\boldsymbol{A}^{+}$ 下对称），唯一性给出结论。
$(\boldsymbol{A}^{\mathrm{T}})^{+} = (\boldsymbol{A}^{+})^{\mathrm{T}}$：$\boldsymbol{A}$ 的条件经转置后恰好是 $\boldsymbol{A}^{\mathrm{T}}$ 的四个条件，且由 $(\boldsymbol{A}\boldsymbol{A}^{+})^{\mathrm{T}} = \boldsymbol{A}\boldsymbol{A}^{+}$ 得 $(\boldsymbol{A}^{+})^{\mathrm{T}}\boldsymbol{A}^{\mathrm{T}}$ 对称，故唯一性给出结论。其余转置/共轭变体同理。
式 $\boldsymbol{A}^{+} = (\boldsymbol{A}^{\mathrm{T}}\boldsymbol{A})^{+}\boldsymbol{A}^{\mathrm{T}}$：由 $\boldsymbol{A}^{\mathrm{T}}\boldsymbol{A}$ 的伪逆结合 SVD（式~\eqref{eq:矩阵:SVD}）代入验证四个条件即可。幂等性：$(\boldsymbol{A}\boldsymbol{A}^{+})^{2} = \boldsymbol{A}\boldsymbol{A}^{+}\boldsymbol{A}\boldsymbol{A}^{+} = \boldsymbol{A}\boldsymbol{A}^{+}$，幂等矩阵的迹等于其秩（见式~\eqref{eq:矩阵:rank恒等式}）。证毕。
\end{proof}

\subsection*{矩阵导数的基本法则}

\begin{proof}
设 $\boldsymbol{A}$、$\boldsymbol{B}$ 为与 $\boldsymbol{X}$ 无关的常数矩阵，由求导的线性性质与乘积法则，
\[
\frac{\partial\operatorname{tr}(\boldsymbol{A}\boldsymbol{X})}{\partial\boldsymbol{X}} = \boldsymbol{A}^{\mathrm{T}}, \qquad
\frac{\partial\operatorname{tr}(\boldsymbol{X}^{\mathrm{T}}\boldsymbol{A}\boldsymbol{X})}{\partial\boldsymbol{X}} = \boldsymbol{A}\boldsymbol{X}+\boldsymbol{A}^{\mathrm{T}}\boldsymbol{X}, \qquad
\frac{\partial\operatorname{tr}(\boldsymbol{A}\boldsymbol{X}\boldsymbol{B})}{\partial\boldsymbol{X}} = \boldsymbol{A}^{\mathrm{T}}\boldsymbol{B}^{\mathrm{T}}.
\]
逐元素验证第一式：$\operatorname{tr}(\boldsymbol{A}\boldsymbol{X}) = \sum_{ij}A_{ij}X_{ji}$，对 $X_{pq}$ 求导得 $A_{qp} = (\boldsymbol{A}^{\mathrm{T}})_{pq}$。第二式先由乘积法则，再由对称性 $\boldsymbol{A}\boldsymbol{X}+\boldsymbol{A}^{\mathrm{T}}\boldsymbol{X}$ 合并。第三式同理。逆矩阵的导数由 $\boldsymbol{X}\boldsymbol{X}^{-1} = \boldsymbol{I}$ 两边求导得 $\dfrac{\partial\boldsymbol{X}^{-1}}{\partial X_{pq}} = -\boldsymbol{X}^{-1}\boldsymbol{J}_{pq}\boldsymbol{X}^{-1}$，用单元素矩阵（式~\eqref{eq:矩阵:jordan块} 前文的记号）表述即为书中所列结果。证毕。
\end{proof}

\section{Vandermonde 行列式的证明}
\label{sec:vandermonde证明}

\begin{proof}
对式~\eqref{eq:矩阵:vandermonde} 用归纳。$n=1$ 时 $\det\boldsymbol{V} = 1$。设 $n-1$ 阶成立。把第 $n$ 行后的行依次减去第 $n-1$ 行的 $v_n$ 倍，最后一行减第 $n-1$ 行的 $v_n$ 倍，沿 $v_n$ 因子提取：
\[
\det\boldsymbol{V}_n = \prod_{i=1}^{n-1}(v_n-v_i)\det\boldsymbol{V}_{n-1} = \prod_{i=1}^{n-1}(v_n-v_i)\prod_{i>j,\,i,j<n}(v_i-v_j) = \prod_{i>j}(v_i-v_j).
\]
证毕。
\end{proof}

\section{Kronecker 积性质的证明}
\label{sec:kronecker证明}

\begin{proof}
按定义（式~\eqref{eq:矩阵:kronecker}）直接分块乘可证混合积性质
\[
(\boldsymbol{A}\otimes\boldsymbol{B})(\boldsymbol{C}\otimes\boldsymbol{D}) = \boldsymbol{A}\boldsymbol{C}\otimes\boldsymbol{B}\boldsymbol{D}.
\]
其转置、逆与秩性质由该式与块状结构直接推出。$\operatorname{tr}(\boldsymbol{A}\otimes\boldsymbol{B}) = \operatorname{tr}(\boldsymbol{A})\operatorname{tr}(\boldsymbol{B})$：按定义求和即得
\[
\operatorname{tr}(\boldsymbol{A}\otimes\boldsymbol{B}) = \sum_{i,k}A_{ii}B_{kk} = \operatorname{tr}(\boldsymbol{A})\operatorname{tr}(\boldsymbol{B}).
\]
特征值性质：设 $\boldsymbol{A}\boldsymbol{x}_i = \lambda_i\boldsymbol{x}_i$、$\boldsymbol{B}\boldsymbol{y}_j = \mu_j\boldsymbol{y}_j$，则
\[
(\boldsymbol{A}\otimes\boldsymbol{B})(\boldsymbol{x}_i\otimes\boldsymbol{y}_j) = \boldsymbol{A}\boldsymbol{x}_i\otimes\boldsymbol{B}\boldsymbol{y}_j = \lambda_i\mu_j\,(\boldsymbol{x}_i\otimes\boldsymbol{y}_j),
\]
故 $\boldsymbol{A}\otimes\boldsymbol{B}$ 的特征值为 $\lambda_i\mu_j$。由特征值行列式关系（式~\eqref{eq:矩阵:特征值迹行列式}）
\[
\det(\boldsymbol{A}\otimes\boldsymbol{B}) = \prod_{i,j}\lambda_i\mu_j = \left(\prod_{i}\lambda_i\right)^{\operatorname{rank}(\boldsymbol{B})}\left(\prod_{j}\mu_j\right)^{\operatorname{rank}(\boldsymbol{A})} = \det(\boldsymbol{A})^{\operatorname{rank}(\boldsymbol{B})}\det(\boldsymbol{B})^{\operatorname{rank}(\boldsymbol{A})}.
\]
vec 算子公式（式~\eqref{eq:矩阵:vec}）$\operatorname{vec}(\boldsymbol{A}\boldsymbol{X}\boldsymbol{B}) = (\boldsymbol{B}^{\mathrm{T}}\otimes\boldsymbol{A})\operatorname{vec}(\boldsymbol{X})$ 可逐元素验证：$(\boldsymbol{A}\boldsymbol{X}\boldsymbol{B})_{ij} = \sum_{kl}A_{ik}X_{kl}B_{lj}$，右端按 vec 的定义将 $X_{kl}$ 排列为 $(\boldsymbol{B}^{\mathrm{T}}\otimes\boldsymbol{A})$ 的行即可。证毕。
\end{proof}

\section{指数矩阵函数性质的证明}
\label{sec:指数证明}

\subsection*{$\mathrm{e}^{A}\mathrm{e}^{B} = \mathrm{e}^{A+B}$（当 $AB = BA$）}

\begin{proof}
设 $\boldsymbol{A}\boldsymbol{B} = \boldsymbol{B}\boldsymbol{A}$，按定义（式~\eqref{eq:矩阵:矩阵指数}）展开
\[
\mathrm{e}^{\boldsymbol{A}}\mathrm{e}^{\boldsymbol{B}} = \sum_{m=0}^{\infty}\sum_{n=0}^{\infty}\frac{\boldsymbol{A}^{m}\boldsymbol{B}^{n}}{m!\,n!}.
\]
由可交换性，含 $\boldsymbol{A}^{p}\boldsymbol{B}^{q}$（$p+q = k$）的各项系数为 $\sum_{m+n=k}\dfrac{1}{m!\,n!}$，利用二项式系数
\[
\sum_{m+n=k}\frac{1}{m!\,n!} = \frac{1}{k!}\binom{k}{m} = \frac{1}{k!},
\]
注意 $\sum_{m=0}^{k}\binom{k}{m} = 2^{k}$ 并不直接适用；正确做法是
\[
\sum_{m=0}^{k}\frac{1}{m!(k-m)!} = \frac{1}{k!}\sum_{m=0}^{k}\binom{k}{m} = \frac{2^{k}}{k!}.
\]
但 $\mathrm{e}^{\boldsymbol{A}+\boldsymbol{B}}$ 的展开系数为 $\dfrac{(\boldsymbol{A}+\boldsymbol{B})^{k}}{k!}$，其中由二项式定理（需 $\boldsymbol{A}\boldsymbol{B} = \boldsymbol{B}\boldsymbol{A}$）$(\boldsymbol{A}+\boldsymbol{B})^{k} = \sum_{m=0}^{k}\binom{k}{m}\boldsymbol{A}^{m}\boldsymbol{B}^{k-m}$，恰好与上述双和一致。故 $\mathrm{e}^{\boldsymbol{A}}\mathrm{e}^{\boldsymbol{B}} = \mathrm{e}^{\boldsymbol{A}+\boldsymbol{B}}$。证毕。
\end{proof}

\subsection*{$\det(\mathrm{e}^{A}) = \mathrm{e}^{\operatorname{tr}(A)}$}

\begin{proof}
先对可对角化情形证明。$\boldsymbol{A} = \boldsymbol{V}\operatorname{diag}(\lambda_i)\boldsymbol{V}^{-1}$ 时
\[
\mathrm{e}^{\boldsymbol{A}} = \boldsymbol{V}\operatorname{diag}(\mathrm{e}^{\lambda_i})\boldsymbol{V}^{-1},
\]
故 $\det(\mathrm{e}^{\boldsymbol{A}}) = \prod_{i}\mathrm{e}^{\lambda_i} = \mathrm{e}^{\sum_i\lambda_i} = \mathrm{e}^{\operatorname{tr}(\boldsymbol{A})}$。对亏损矩阵，用 Jordan 标准形（式~\eqref{eq:矩阵:jordan标准形}）：$\boldsymbol{A} = \boldsymbol{P}\boldsymbol{J}\boldsymbol{P}^{-1}$，则 $\mathrm{e}^{\boldsymbol{A}} = \boldsymbol{P}\mathrm{e}^{\boldsymbol{J}}\boldsymbol{P}^{-1}$。Jordan 块的上三角因子给出 $\det(\mathrm{e}^{\boldsymbol{J}}) = \prod_{k}\mathrm{e}^{\lambda_k}$（三角阵行列式为对角元之积，上三角部分无贡献），仍得 $\det(\mathrm{e}^{\boldsymbol{A}}) = \mathrm{e}^{\operatorname{tr}(\boldsymbol{A})}$。证毕。
\end{proof}

\subsection*{$\mathrm{d}/\mathrm{d}t\,\mathrm{e}^{tA} = A\mathrm{e}^{tA}$}

\begin{proof}
按定义对级数逐项求导：
\[
\frac{\mathrm{d}}{\mathrm{d}t}\mathrm{e}^{t\boldsymbol{A}} = \frac{\mathrm{d}}{\mathrm{d}t}\sum_{n=0}^{\infty}\frac{t^{n}\boldsymbol{A}^{n}}{n!} = \sum_{n=1}^{\infty}\frac{t^{n-1}\boldsymbol{A}^{n}}{(n-1)!} = \boldsymbol{A}\sum_{n=1}^{\infty}\frac{t^{n-1}\boldsymbol{A}^{n-1}}{(n-1)!} = \boldsymbol{A}\mathrm{e}^{t\boldsymbol{A}}.
\]
由幂级数在有限区间内一致收敛可交换求导与求和。$\mathrm{e}^{t\boldsymbol{A}}$ 与 $\boldsymbol{A}$ 可交换，故亦等于 $\mathrm{e}^{t\boldsymbol{A}}\boldsymbol{A}$。证毕。
\end{proof}
